\subsection{Structural pushforward and twisted higher Chow groups}\label{sec:global-chow}
This subsection prepares the cycle-theoretic input for Section~\ref{sec:global-spectrum}.
We relate twisted higher Chow groups to ordinary ones through the pushforward along the structure map of a Severi--Brauer variety: this pushforward factors through the twisted higher Chow groups, the factored map is an isomorphism on the diagonal and surjective on the next row, and composing it with the reduced cycle map gives the usual pushforward.
After $\ell$-localization, the reduced cycle map is an isomorphism on both rows in weights $r\ge2$, except for $\ell=2$ over number fields with real places (Proposition~\ref{prop:global-two-chow-rows}); for $\CH^2(F,3;D)$ this exception is removed in Proposition~\ref{prop:global-real-weight-two}.
These isomorphisms are used for $K_3(D)$ and for the group $\pi_4$ of the comparison fiber.

Let $F$ be a global field and $D$ a central division algebra over $F$.
At a closed support $T$, the reduced rank identifies $K_0(D_{\kappa(T)})$ with $\ind(D_{\kappa(T)})\Z$, inducing the reduced cycle map $\rho$ of Definition~\ref{def:paper-twisted-chow}.
We compare $\rho$ with the \emph{structural trace}, the pushforward along the structure map of a Severi--Brauer variety.
Write $\mathcal F^rK(D)$ for the homotopy coniveau tower.
The homotopy group $\pi_m$ of the weight-$r$ layer is $\CH^r(F,m;D)$, and $\pi_i\mathcal F^rK(D)=0$ for $i<r$.
The weight-zero layer has only the rank group in degree zero \cite[Corollary 6.1.4 and Section 6.6]{KahnLevine2010}.
Over a global function field, which is imperfect, the identification of the layers is available by Lemma~\ref{lem:app-twisted-tower}.

\subsubsection{Structural pushforward and reduced rank}
Our first goal is Lemma~\ref{lem:integral-structural-push}, which expresses the twisted higher Chow groups on the diagonal and the next row through higher Chow groups of a Severi--Brauer variety, compatibly with the reduced cycle map.
Its proof uses the following boundary correction.

\begin{lemma}\label{lem:paper-structural-boundary-correction}
Let $F$ be an infinite field, let $B$ be a central simple algebra over $F$, put $X=\SB(B)$ and $n=\dim X$, let $r\ge1$, and write $f_*\colon z^{n+r}(X,\bullet)\to z^r(F,\bullet)$ for the proper pushforward along the structure map $f\colon X\to\Spec F$.
If a chain $c$ of degree $r$ satisfies $f_*c=0$, then there is a chain $W$ of degree $r+1$ with
\[
 f_*W=0,\qquad\partial W=c.
\]
\end{lemma}
\begin{proof}
The chain $c$ has codimension $n+r=\dim(X\times\Delta^r)$, so its components are closed points, and proper intersection with the faces forces their images in $\Delta^r$ to avoid all faces.
Over each interior closed point $q\in\Delta^r$, the component of $c$ is a zero-cycle of degree zero on $X_{\kappa(q)}$.
The degree map on $\CH_0$ of a Severi--Brauer variety is injective \cite[Corollary 7.3]{ChernousovMerkurjev2006}, so the component is a sum of pushforwards of divisors of rational functions on proper normal integral curves $C\to X_{\kappa(q)}$.
By normality of $C$, a rational function $\psi$ with $\operatorname{div}(\psi)\ne0$ defines a nonconstant morphism $\psi\colon C\to\PP^1_{\kappa(q)}$.
This morphism of proper curves is finite, and it is flat because $\psi_*\mathcal O_C$ is torsion-free over the regular curve $\PP^1_{\kappa(q)}$.

Set $L=\kappa(q)$ and place the corresponding $L$-point $a=(a_0,\ldots,a_{r+1})$ on the face $t_0=0$ of $\Delta^{r+1}_L$.
Thus $a_0=0$ and $a_j\ne0$ for $j>0$.
We choose $b_j\in L$ with $\sum_j b_j=0$, $b_j\ne0$, and $a_i b_j-a_j b_i\ne0$ for $i\ne j$; this avoids finitely many proper linear subspaces of the hyperplane $\sum_jb_j=0$ over the infinite field $L$.
The line $t_j=a_j+b_j u$, parametrized by $u\in\mathbb A^1_L$, meets the codimension-one faces transversely at distinct points $u_j=-a_j/b_j$ and avoids all faces of higher codimension.
We interpolate a polynomial $h$ with $h(u_0)=1$ and $h(u_j)=0$ for $j>0$ and set $\varphi(u)=[1-h(u):h(u)]$.
Pullback along $\varphi$ and along the constant map with value $\infty$ gives
\[
 C\times_{\PP^1_L,\varphi}\mathbb A^1_L,\qquad
 C\times_{\PP^1_L,\infty}\mathbb A^1_L,
\]
both finite flat of the same degree over the line, where $0=[0:1]$ and $\infty=[1:0]$ in $\PP^1_L$.
The difference of their fundamental cycles gives an admissible chain $W_\psi$ by proper pushforward to $X_L\times\Delta^{r+1}_L$ followed by finite pushforward to $F$.
Equal degrees give $f_*W_\psi=0$.
By transversality and finite flatness, restriction to $t_j=0$ gives the difference of the fiber cycles at $u_j$, namely $\psi^{-1}(0)-\psi^{-1}(\infty)$ for $j=0$ and zero for $j>0$.
The pushforward has no extra degree factor since $L=\kappa(q)$, and the boundary sign of $t_0=0$ is $+1$.
Summing the chains $W_\psi$ over the divisors, we obtain a chain $W$ with
\[
 f_*W=0,\qquad\partial W=c.
\]
\end{proof}

\begin{lemma}\label{lem:integral-structural-push}
Let $F$ be an infinite field, let $D$ be a central division algebra over $F$, let $B$ be a central simple algebra over $F$ with $[B]=[D]$, and put $X=\SB(B)$ and $n=\dim X$.
For every $r\ge1$ the structure map $f\colon X\to\Spec F$ induces
\[
 \CH^{n+r}(X,r+1)\twoheadrightarrow\CH^r(F,r+1;D),\qquad
 \CH^{n+r}(X,r)\xrightarrow{\sim}\CH^r(F,r;D).
\]
The composite with $\rho$ is the usual proper pushforward along the structure map.
These maps commute with restriction and transfer for finite extensions.
\end{lemma}

\begin{proof}
We first show that $f_*$ factors through the twisted cycle complex.
In each simplicial degree, let $W\subset X\times\Delta^m$ be an integral cycle with image $T\subset\Delta^m$.
Then
\[
 f_*[W]=
 \begin{cases}
  [\kappa(W):\kappa(T)]\,[T], & \dim T=\dim W,\\
  0, & \dim T<\dim W.
 \end{cases}
\]
In the first case $X$ has a $\kappa(W)$-rational point, so $[\kappa(W):\kappa(T)]$ is divisible by $\ind(D_{\kappa(T)})$.
Consequently $f_*$ factors through the twisted cycle complex $z^r(F,\bullet;D)\subset z^r(F,\bullet)$ embedded by reduced rank, compatibly with face restrictions, flat base change, and proper transfer.

We next show surjectivity on chains.
For $m=r,r+1$, an admissible support $T$ is a closed point or an integral curve, respectively.
To lift the coefficient generator $\ind(D_{\kappa(T)})[T]$, we take the closure of a closed point of $X_{\kappa(T)}$ of degree $\ind(D_{\kappa(T)})$.
When $T$ is a curve, the closure is proper over $T$, and its closed fibers are proper closed subsets of the closure, an integral curve; hence they are zero-dimensional, and the closure is finite over $T$.
It inherits the proper intersection conditions with faces from $T$.
This proves surjectivity on chains in both degrees.

We lift a cycle in degree $r+1$ of the twisted cycle complex to a chain $Z$; then $f_*\partial Z=0$.
For the chain $W$ of Lemma~\ref{lem:paper-structural-boundary-correction} with $c=\partial Z$, the difference $Z-W$ is a cycle lifting the given cycle.
In degree $r$ the chain group one degree lower vanishes, giving surjectivity.
Suppose that a degree-$r$ chain $Y$ and a chain $V$ satisfy $f_*Y=\partial V$.
We choose a lift $\widetilde V$ of $V$ and subtract $\partial\widetilde V$ from $Y$.
Then $f_*(Y-\partial\widetilde V)=0$, so $Y-\partial\widetilde V$ is also a boundary by Lemma~\ref{lem:paper-structural-boundary-correction} with $c=Y-\partial\widetilde V$.
This proves injectivity in degree $r$.
\end{proof}

\subsubsection{Reduced norms on the diagonal and the next row}
Next we use Lemma~\ref{lem:integral-structural-push} to study the reduced cycle map $\rho$ on $\CH^r(F,m;D)$ for $m=r,r+1$.

\begin{proposition}\label{prop:global-two-chow-rows}
Let $F$ be a global field and $\ell$ a prime.
If $F$ is a number field, we assume that $\ell$ is odd or that $F$ has no real places.
For every $r\ge2$,
\[
 \rho\colon\CH^r(F,m;D)_{(\ell)}
    \xrightarrow{\sim}\CH^r(F,m)_{(\ell)}
    \qquad(m=r,r+1).
\]
The isomorphisms commute with restriction and transfer.
\end{proposition}
\begin{proof}
First let $\ell=\operatorname{char}F$, so that $[F:F^\ell]=\ell$.
We split the prime-to-$\ell$ component of the Brauer class over an extension $F'/F$ of degree prime to $\ell$.
Since $\cor_{F'/F}\res_{F'/F}=[F':F]$, the kernel and cokernel of $\rho_{(\ell)}$ are detected after extension.
There we apply the rank-$1$ Frobenius comparison of Proposition~\ref{prop:app-frobenius}.

Now let $\ell\ne\operatorname{char}F$, and put $X=\SB(D)$, $n=\dim X$, and $R=\Z_{(\ell)}$.
We consider the fiber of the structural trace
\[
 \mathcal C_r=\fib\!\left[
 R\Gamma_{\mot}(X,R(n+r))[2n]
   \xrightarrow{f_*}R\Gamma_{\mot}(F,R(r))\right].
\]
Since $\CH^r(F,m)=H^{2r-m}_{\mot}(F,\Z(r))$ and $\CH^{n+r}(X,m)=H^{2r-m}\bigl(R\Gamma_{\mot}(X,\Z(n+r))[2n]\bigr)$, the rows $m=r$ and $m=r+1$ correspond to degrees $r$ and $r-1$ of the trace.
By the hypothesis on $\ell$, $\cd_\ell F\le2$.
At a point of dimension $a$ and codimension $c=n-a$ in $X$, purity gives weight $a+r$, and the residue field has $\ell$-cohomological dimension at most $a+2\le a+r$ (as $r\ge2$).
The layers of the support filtration therefore lie in the range of the Beilinson--Lichtenbaum comparison \cite{Voevodsky}, \cite[Theorem 2.3(2)]{Yamazaki2013}.
Since the cycle class maps commute with proper trace, this identifies the finite-coefficient trace with the \'etale trace in the required degrees.

Over $F^{\mathrm{sep}}$, $X$ is $\PP^n$.
The trace shifted by $[2n]$ is an isomorphism on the top geometric group $\Z/\ell(r)$.
The remaining groups lie in geometric cohomological degrees $-2,-4,\ldots,-2n$.
Modulo $\ell$, Hochschild--Serre gives
\[
 H^q(\mathcal C_r/\ell)=0\qquad(q\ge1).
\]
With rational coefficients, we descend the projective bundle decomposition from a finite splitting field by restriction and transfer, using the hyperplane class $c_1(\omega_X^{-1})/(n+1)$.
The trace is projection onto the top component, so
\[
 (\mathcal C_r)_\Q\simeq
 \bigoplus_{k=1}^{n}R\Gamma_{\mot}(F,\Q(r+k))[2k].
\]
For number fields the weight decomposition of Borel's calculation \cite{Borel} via the Chern character gives $H^i_{\mot}(F,\Q(b))=0$ for $b\ge2$ and $i\ne1$ (see the proof of \cite[Theorem 3.2]{Scholbach2011}).
For global function fields the groups in this range vanish by Harder's theorem \cite{Harder}; this is consistent with localization, since the positive $K$-groups of the finite residue fields are torsion \cite{QuillenFinite}.
Thus $H^q((\mathcal C_r)_\Q)=0$ for $q\ge0$.

Combining the two computations, the coefficient sequence shows that multiplication by $\ell$ is invertible on $H^q(\mathcal C_r)$ for $q\ge2$.
These groups are $\ell$-primary torsion and hence zero.
The group $H^1(\mathcal C_r)$ is $\ell$-divisible.
Consequently the structural trace is an isomorphism in degrees $q\ge2$ and is surjective with $\ell$-divisible kernel in degree one.
Write $\widetilde f_*$ for the maps $\CH^{n+r}(X,m)\to\CH^r(F,m;D)$ of Lemma~\ref{lem:integral-structural-push}, so that $f_*=\rho\circ\widetilde f_*$.
On the diagonal, $\widetilde f_*$ is an isomorphism and identifies $\rho$ with the trace in degree $r$.
On the next row, $f_*$ and $\widetilde f_*$ are both surjective, so $\rho$ is surjective and $\ker\rho=\widetilde f_*(\ker f_*)$.
This kernel vanishes for $r\ge3$, since $f_*$ is injective in degree $r-1\ge2$, and it is $\ell$-divisible for $(r,m)=(2,3)$.
In the case $(r,m)=(2,3)$, restriction and transfer show that $\ker\rho$ is killed by the degree of a finite splitting field.
The prime-to-$\ell$ part of this degree is a unit, so this $\ell$-divisible group is killed by a power of $\ell$ and is therefore zero.
\end{proof}

\subsubsection{The trace relative to the real places}
Proposition~\ref{prop:global-two-chow-rows} excludes the prime two over number fields with real places.
For $\CH^2(F,3;D)$ we remove this restriction in Proposition~\ref{prop:global-real-weight-two}, using the structural trace relative to the real places.

\begin{proposition}\label{prop:global-real-weight-two}
For a number field $F$ with a real place and any central division algebra $D$ over $F$,
\[
 \rho\colon\CH^2(F,3;D)_{(2)}\xrightarrow{\sim}\CH^2(F,3)_{(2)}.
\]
\end{proposition}
\begin{proof}
Set $X=\SB(D)$, $n=\dim X$, and $R=\Z_{(2)}$.
The structural trace in degree one is
\[
 f_*\colon H^{2n+1}_{\mot}(X,R(n+2))\longrightarrow H^1_{\mot}(F,R(2)),
\]
and the target of $\rho$ is $\CH^2(F,3)_{(2)}=K_3^{\mathrm{ind}}(F)_{(2)}$.

We write $R\Gamma_{\mot,\mathrm{rel}}$ and $H^i_{\mathrm{rel}}$ for the fiber of restriction to the algebraic real closures $F_v^{\mathrm{rc}}$ at all real places $v$, and its cohomology.
For $M=R,\Q,\Q/R$, set
\[
 \mathcal C_{\mathrm{rel}}(F,M)=\fib\left(
 R\Gamma_{\mot,\mathrm{rel}}(X,M(n+2))[2n]
 \xrightarrow{f_*} R\Gamma_{\mot,\mathrm{rel}}(F,M(2))
 \right).
\]
We first show that $H^0\mathcal C_{\mathrm{rel}}(F,\Q/R)$ and $H^1\mathcal C_{\mathrm{rel}}(F,\Q/R)$ vanish.

We apply finite-coefficient motivic--\'etale comparison to the support filtration.
At a point of dimension $a$ on $X$, purity shifts the weight to $a+2$.
Without real relativization the virtual $2$-cohomological dimension is at most $a+2$.
In the relative calculation, the constant field in the residue field is a finite number field, and its real and complex branches account for the restriction targets.
By Jannsen \cite[Proposition 1.2 and Lemma 1.3]{Jannsen}, restriction to the real branches is surjective in degree $a+2$ and an isomorphism above that degree.
The complex branches have cohomological dimension at most $a$.
Thus the comparison cofiber above the Beilinson--Lichtenbaum range \cite{Voevodsky} vanishes after taking the relative fiber.
For $x\in X$, let $k_x$ be the algebraic closure of $F$ in $\kappa(x)$; the residue fields on the real-closure side are the factors of $\kappa(x)\otimes_FF_v^{\mathrm{rc}}$.
Closed supports over a real closure descend to finite intermediate fields and lie in the base change of their finite images; since codimension is preserved, base-changed supports are cofinal, so the relative layers can be computed residue field by residue field \cite[Proposition 1.2, Lemma 1.3(b), Remark 1.4(a)]{Jannsen}.
Taking the direct limit over finite coefficients, we compute $\mathcal C_{\mathrm{rel}}(F,\Q/R)$ by \'etale cohomology.

Over $F^{\mathrm{sep}}$ the trace is an isomorphism on $(\Q/R)(2)$ in top geometric degree $2n$, so after shifting by $[2n]$ the geometric fiber has cohomology only in even degrees at most $-2$, with $(\Q/R)(3)$ in degree $-2$ for $n\ge1$.
The relative cohomological dimension of the base is at most two, so positive degrees vanish and the only possible contribution in degree zero is $H^2_{\mathrm{rel}}(F,(\Q/R)(3))$.
The motivic coefficient triangle gives
\[
 H^2(F,(\Q/R)(3))\simeq K_3^M(F)\otimes R.
\]
Indeed, the rational groups $H^2_{\mot}(F,\Q(3))$ and $H^3_{\mot}(F,\Q(3))$ on either side of the $\Q/R$-coefficient term vanish by $K_4(F)_\Q=0$ \cite{Borel} and $K_3^M(F)_\Q=0$, respectively.
The same holds at algebraic real closures by taking the direct limit over finite number fields.
Using the description of $K_3^M(F)$ by signs at the real places \cite{BassTate}, we have
\[
 H^2(F,(\Q/R)(3))\simeq\prod_{v\,\mathrm{real}}H^2(F_v^{\mathrm{rc}},(\Q/R)(3)).
\]
Complex conjugation acts on $(\Q/R)(3)$ by $-1$, so
\[
 H^1\bigl(F_v^{\mathrm{rc}},(\Q/R)(3)\bigr)=(\Q/R)(3)/2(\Q/R)(3)=0.
\]
The relative exact sequence gives $H^2_{\mathrm{rel}}(F,(\Q/R)(3))=0$.
Hence $H^i\mathcal C_{\mathrm{rel}}(F,\Q/R)=0$ for $i\ge0$.

Next we compute with rational coefficients.
Using the hyperplane class above, we descend the rational projective bundle decomposition.
Since the structural trace projects onto the top component,
\[
 \mathcal C_{\mathrm{rel}}(F,\Q)\simeq
 \bigoplus_{j=3}^{n+2}
 R\Gamma_{\mot,\mathrm{rel}}(F,\Q(j))[2j-4].
\]
For $i\ge1$ the absolute group $H^{2j+i-4}_{\mot}(F,\Q(j))$ vanishes since the degree exceeds the weight except for $j=3$ and $i=1$, where it is $K_3^M(F)_\Q=0$.
The preceding real-closure group $H^{2j+i-5}_{\mot}(F_v^{\mathrm{rc}},\Q(j))$ is handled in the same way.
The exceptions within the weight bound are $H^2_{\mot}(-,\Q(3))$, $K_3^M(-)_\Q$, and $K_4^M(-)_\Q$, all zero by the number field calculation and filtered colimits.
Therefore $H^i\mathcal C_{\mathrm{rel}}(F,\Q)=0$ for $i\ge1$, and the coefficient triangle gives $H^i\mathcal C_{\mathrm{rel}}(F,R)=0$ for $i\ge1$.

We now return to the absolute trace.
For a field extension $k/F$ and $M$ as above, write $X_k=X\times_Fk$ and define
\[
 \mathcal C_{\mathrm{abs}}(k,M)=\fib\left(
 R\Gamma_{\mot}(X_k,M(n+2))[2n]
 \xrightarrow{f_*} R\Gamma_{\mot}(k,M(2))
 \right).
\]
The exact sequence with the absolute trace fiber gives
\[
 H^1\mathcal C_{\mathrm{abs}}(F,R)
 \xrightarrow{\sim}\prod_{v\,\mathrm{real}}H^1\mathcal C_{\mathrm{abs}}(F_v^{\mathrm{rc}},R),
 \qquad
 H^2\mathcal C_{\mathrm{abs}}(F,R)
 \hookrightarrow\prod_{v\,\mathrm{real}}H^2\mathcal C_{\mathrm{abs}}(F_v^{\mathrm{rc}},R).
\]
By Lemma~\ref{lem:integral-structural-push} with $r=2$, the structure map induces a surjection $\CH^{n+2}(X,3)\twoheadrightarrow\CH^2(F,3;D)$ whose composite with $\rho$ is $f_*$.
The same holds over $F_v^{\mathrm{rc}}$ and over finite extensions of $F$, and these surjections commute with restriction and transfer.
Over a real closure, the index is $1$ or $2$; after passing by Morita equivalence to degree $2$ in the second case, $\rho$ is an isomorphism in weight two \cite[Theorems 6.5.3 and 8.1.4(1)]{KahnLevine2010}.
Hence $f_*$ is surjective onto $H^1_{\mot}(F_v^{\mathrm{rc}},R(2))$, so the connecting map $H^1_{\mot}(F_v^{\mathrm{rc}},R(2))\to H^2\mathcal C_{\mathrm{abs}}(F_v^{\mathrm{rc}},R)$ vanishes.
By naturality and the injectivity in $H^2$ displayed above, the connecting map $H^1_{\mot}(F,R(2))\to H^2\mathcal C_{\mathrm{abs}}(F,R)$ also vanishes.
Thus $f_*$ is surjective over $F$, and since $f_*$ factors through $\rho$, so is $\rho$ onto $K_3^{\mathrm{ind}}(F)_{(2)}$.
For injectivity, lift a kernel element in weight two along the surjection of Lemma~\ref{lem:integral-structural-push} to $H^{2n+1}_{\mot}(X,R(n+2))$ and then to an element $y$ of $H^1\mathcal C_{\mathrm{abs}}(F,R)$.
Choose $k=F(\alpha)$ where the minimal polynomial of $\alpha$ is a monic polynomial of degree $\deg D$ chosen by weak approximation: at each finite place where $D$ is ramified, close to an irreducible polynomial defining the unramified extension of degree $\deg D$ (Krasner's lemma), and at each real place, close to a polynomial with distinct real roots (simple roots form an open condition); then all real places split completely in $k$, and the local degrees at the finite ramified places are divisible by the index.
If no finite place is ramified, then $\ind D\le2$ and $k=F$ suffices.
Brauer--Hasse--Noether gives $\ind(D_k)\le2$.
Place the component of $y$ at each real place on one real branch of $k$ and use zero on the others.
Since the relative computation applies over $k$ to $X_k=\SB(D_k)$, the isomorphism in $H^1$ onto the real components, applied over $k$, lifts the resulting tuple to $y_k\in H^1\mathcal C_{\mathrm{abs}}(k,R)$.
At each real place $v$ of $F$, corestriction adds the components over the real places of $k$ above $v$, so $\cor_{k/F}y_k$ and $y$ have the same real components, and the injectivity in $H^1$ gives $\cor_{k/F}y_k=y$.
Since $\rho$ is an isomorphism at index at most two, after passing by Morita equivalence to degree $2$ at index two \cite[Theorems 6.5.3 and 8.1.4(1)]{KahnLevine2010}, $y_k$ has zero image in the twisted weight-two group.
Transferring back shows that the original kernel element is zero.
\end{proof}
